
\subsubsection{2.1 Spurious Correlations and Worst-Group Accuracy}

Spurious correlations --- statistical associations between input features and labels that hold in training data but do not generalize --- are a well-documented failure mode of empirical risk minimization \citep{sagawa2019distributionally}. On the Waterbirds benchmark \citep{sagawa2019distributionally}, ResNet-50 models trained with ERM exploit background (land vs. water) as a shortcut for bird species classification, achieving high average accuracy but failing on minority groups (e.g., landbirds on water). WGA --- accuracy on the worst-performing group --- is the standard measure for spurious correlation robustness.

Multiple training-time interventions improve WGA over ERM: GroupDRO \citep{sagawa2019distributionally} reweights groups dynamically; JTT \citep{Liu2021} upweights misclassified examples; SAM \citep{foret2021sharpness} improves sharpness-aware generalization. These methods achieve WGA improvements but the mechanisms by which they do so --- specifically whether improvement is mediated by backbone representation change or head calibration --- remain underspecified. This study addresses the mechanistic question for GroupDRO and DFR.

\subsubsection{2.2 Deep Feature Reweighting and Last-Layer Retraining}

Kirichenko et al. [2022] demonstrated that much WGA improvement can be achieved by retraining only the final classification head on group-balanced held-out data while freezing the backbone (DFR). Subsequent work [LaBonte et al., 2023; Hill et al., 2025] refined this understanding: the key condition is group balance in the held-out set used for head retraining. DFR and ERM backbones are identical by design; Izmailov et al. [2022] release both checkpoints without explicitly quantifying this backbone identity empirically at matching seeds. H-P0 fills this gap (cosine similarity = 1.000000 $\pm$ 1e-14, all 3 seed pairs).

Le et al. [2023] provide a cautionary finding: in medical imaging domains, DFR is insufficient when backbone features strongly encode spurious attributes and backbone-level intervention is required. This is consistent with the framework advanced here: the sufficiency of head-only retraining depends on the backbone's spurious encoding level.

\subsubsection{2.3 Feature Analysis of Robustification Methods}

Izmailov et al. [2022] use a spurious-DFR proxy (s-DFR) to estimate spurious information encoded in ERM backbones (~85\% probe accuracy), but do not report per-method, per-seed probe accuracy with statistical tests comparing GroupDRO to ERM. Murotkar et al. [2024] employ sklearn L-BFGS linear probes on frozen ResNet-50 layer4 features for spurious feature disentanglement --- establishing the probe methodology adopted here.

Park et al. [2025] introduce SCER, which explicitly regularizes the spurious feature subspace during training to reduce background decodability and thereby improve WGA. SCER demonstrates a causal link between reducing background decodability and improving WGA; the present study validates this link from the opposite direction: GroupDRO implicitly achieves what SCER does explicitly. Raymond et al. [2026, preprint] provide corroborating evidence that GroupDRO reshapes backbone representations across all layers, consistent with the H-M2 weight-difference findings; the H-M2 result stands on its own pre-registered evidence independent of this citation.

\subsubsection{2.4 Linear Probing for Representation Analysis}

Linear probing is a standard tool for analyzing backbone representation content \citep{alain2016understanding}. The conventions adopted here (sklearn L-BFGS, C=1e9, frozen layer4 features with global average pooling) match those of Kirichenko et al. [2022] and Murotkar et al. [2024]. The contribution of this study is the per-method, per-seed measurement of backbone spurious encoding with pre-registered statistical tests --- enabling the backbone-vs-head typology.

